\documentclass[11pt,a4paper]{article}
\usepackage[utf8]{inputenc}\usepackage[T1]{fontenc}\usepackage[italian,english]{babel}
\usepackage{amsmath,amssymb,amsthm}\usepackage[margin=2.2cm]{geometry}\usepackage{longtable,booktabs,array,calc}
\usepackage{listings,xcolor}\usepackage{hyperref}\usepackage{url}
\providecommand{\tightlist}{\setlength{\itemsep}{0pt}\setlength{\parskip}{0pt}}
\lstset{basicstyle=\ttfamily\scriptsize,breaklines=true,frame=single,captionpos=b,columns=fullflexible,keepspaces=true,inputencoding=utf8,extendedchars=true}
\title{Proof Pursuit --- Problem 4, part 3\\ \large prove, decisioni degli agenti, codice verificato, letteratura}
\author{Team Proof Pursuit (Researcher: T. Tumini; Referee: gabundos; Creative: M. Renzi; gio3r3gio)}
\date{\today}
\begin{document}\maketitle\tableofcontents
\section*{Come leggere questo rapporto}
Per ogni cella: la richiesta ufficiale, lo stato dichiarato dal team (mai ``risolto'' senza revisione umana), la consegna
con la prova, poi la traccia delle decisioni degli agenti (Researcher: famiglia e motivo; Referee: verdetto ed errore
fatale; Creative: quando e perché è intervenuto) e il codice su cui il tentativo poggia con l'esito della riesecuzione
fidata. DIMOSTRATO, CITATO, VERIFICATO SU INTERVALLO FINITO e CONGETTURA restano distinti. L'architettura del sistema
è descritta nel README del repository.
\section{Problem 4 — Disjoint congruence classes and large gcd}
\subsection{Cella 3 (3 punti) — stato: revisionato}
\subparagraph{\texorpdfstring{Parte 3 (C3) --- Every \(k\) up to
8}{Parte 3 (C3) --- Every k up to 8}}\label{parte-3-c3-every-k-up-to-8}

\textbf{Punteggio:} 3 points · \textbf{Valutazione:} Judged

A first range of sizes: after C1 and C2, the sizes \(k = 5, 6, 7\) and
\(8\) remain. Prove the statement for every \(k \le 8\).

\textbf{Enunciato protetto dato agli agenti (state.json).} \texttt{Prove the statement for every k \textless{}= 8 (sizes 5,6,7,8 remain after cells 1-2). A computation is admissible only with the full exhaustiveness certificate (finite set + proof, pruning rules proved, node counts, wall clock, rerun, decisions for every survivor).}

\textbf{Evidenza registrata in STATUS.md.} loop p4\_c3 attempt\_001: Referee READY\_FOR\_HUMAN (A PASS, B PASS); approvata da Thomas Tumini (human); submission/parte-3.en.md

\subsubsection{Consegna (prova)}
\paragraph{Problem 4 --- Part 3 ---
submission}\label{problem-4-part-3-submission}

\textbf{Declared status: SOLVED (approved).} Complete proof; automatic
Referee READY\_FOR\_HUMAN (judge A, mathematics: PASS; judge B,
evidence: PASS; code re-run by the orchestrator); human approval
recorded in \texttt{runs/p4\_c3/approval.json}.

\subparagraph{1. Result and scope}\label{result-and-scope}

Complete proof of the cell-3 statement for every k in \{3,\ldots,8\}:
(i) Lemma 1 reduces any size-k counterexample (all gcds ≤ 7) to a family
with moduli dividing 420, identical pairwise gcds, still pairwise
disjoint --- full proof given; (ii) exhaustive enumeration of all
k-cliques of the compatibility graph on the 1343 reduced classes finds
no clique for k=3..8, with exhaustiveness certificate: finite set +
proof (Corollary 2), no pruning beyond the definition, node counts (7118
/ 20241 / 56208 / 234719 / 806717 / 3523544), wall clock (2.5 s total,
Apple M4, Python 3.14.7), rerun reproducing all counts exactly, zero
survivors; (iii) the unmodified machinery with threshold k recovers the
trivial partition (non-vacuity); (iv) an independent second
implementation agrees; (v) hand proofs for k=3 and k=4 included so the
cell does not depend on cells 1--2.

\subparagraph{2. Proof}\label{proof}

\subparagraph{Statement proved}\label{statement-proved}

\textbf{Theorem.} Let \(3 \le k \le 8\) and let
\(a_1 \pmod{m_1},\dots,a_k \pmod{m_k}\) be pairwise disjoint congruence
classes (\(m_i\ge 1\)). Then \(\gcd(m_i,m_j)\ge k\) for some \(i<j\).

Throughout, a \emph{counterexample of size \(k\)} is a pairwise disjoint
family of \(k\) classes with \(\gcd(m_i,m_j)\le k-1\) for all \(i<j\).
We use only the definitions and the criterion \((*)\) from the problem
statement: two classes meet iff \(\gcd(m_i,m_j)\mid a_i-a_j\).

\subparagraph{Step 0. Two preliminary
facts}\label{step-0.-two-preliminary-facts}

\textbf{(P1)} In a pairwise disjoint family with \(k\ge2\), every
\(m_i\ge 2\). Indeed, if \(m_i=1\) then \(\gcd(m_i,m_j)=1\mid a_i-a_j\)
for any \(j\neq i\), so the classes \(i,j\) meet by \((*)\).

\textbf{(P2)} Every integer \(g\) with \(1\le g\le 7\) divides
\(N:=420=2^2\cdot3\cdot5\cdot7\). (Check:
\(420=1\cdot420=2\cdot210=3\cdot140=4\cdot105=5\cdot84=6\cdot70=7\cdot60\).)

\subparagraph{Step 1. Reduction lemma (the finite
set)}\label{step-1.-reduction-lemma-the-finite-set}

\textbf{Lemma 1.} Let \(a_1 \pmod{m_1},\dots,a_k \pmod{m_k}\)
(\(k\ge2\)) be pairwise disjoint with \(g_{ij}:=\gcd(m_i,m_j)\le 7\) for
all \(i<j\). Put \(m_i':=\gcd(m_i,N)\) with \(N=420\) and
\(a_i':=a_i \bmod m_i'\) (so \(0\le a_i'<m_i'\)). Then:

\begin{enumerate}
\def\labelenumi{\arabic{enumi}.}
\tightlist
\item
  \(m_i' \mid 420\) and \(m_i'\ge 2\);
\item
  \(\gcd(m_i',m_j')=g_{ij}\) for all \(i<j\);
\item
  the classes \(a_1' \pmod{m_1'},\dots,a_k'\pmod{m_k'}\) are pairwise
  disjoint;
\item
  the \(k\) pairs \((a_i',m_i')\) are pairwise distinct.
\end{enumerate}

\emph{Proof.} (2):
\(\gcd(m_i',m_j')=\gcd(\gcd(m_i,N),\gcd(m_j,N))=\gcd(m_i,m_j,N)=\gcd(g_{ij},N)\).
Since \(1\le g_{ij}\le 7\), (P2) gives \(g_{ij}\mid N\), hence
\(\gcd(g_{ij},N)=g_{ij}\).

(1): \(m_i'\mid N\) by definition. For any \(j\ne i\),
\(g_{ij}\mid m_i\) and \(g_{ij}\mid N\), so \(g_{ij}\mid m_i'\); by (P1)
applied to the original family, \(g_{ij}\ge 2\) --- more precisely,
disjointness of classes \(i,j\) and \((*)\) give
\(g_{ij}\nmid a_i-a_j\), which is impossible for \(g_{ij}=1\); hence
\(g_{ij}\ge2\) and \(m_i'\ge g_{ij}\ge2\).

(3): Fix \(i<j\). By (2), \(\gcd(m_i',m_j')=g_{ij}\). Since
\(g_{ij}\mid m_i'\) and \(g_{ij}\mid m_j'\) (shown in (1)), and
\(a_i'\equiv a_i \pmod{m_i'}\), \(a_j'\equiv a_j\pmod{m_j'}\), we get
\(a_i'-a_j'\equiv a_i-a_j \pmod{g_{ij}}\). The original classes are
disjoint, so by \((*)\) \(g_{ij}\nmid a_i-a_j\), hence
\(g_{ij}\nmid a_i'-a_j'\), and by \((*)\) again the reduced classes
\(i,j\) are disjoint.

(4): Two equal pairs would be the same nonempty class, contradicting
(3). \(\square\)

\textbf{Corollary 2.} If a counterexample of size \(k\) exists with
\(3\le k\le8\), then there exist \(k\) pairwise distinct pairs
\((a_i,m_i)\) with \(m_i\mid 420\), \(m_i\ge2\), \(0\le a_i<m_i\), such
that for every \(i<j\): \(\gcd(m_i,m_j)\le k-1\) and
\(\gcd(m_i,m_j)\nmid a_i-a_j\).

\emph{Proof.} In a counterexample of size \(k\le 8\) all
\(g_{ij}\le k-1\le 7\), so Lemma 1 applies; its conclusions (1)--(4)
together with the preserved gcds (2) give exactly the listed properties.
\(\square\)

\subparagraph{Step 2. The finite search and its
certificate}\label{step-2.-the-finite-search-and-its-certificate}

\textbf{Finite set (certificate item 1).} Let
\(V=\{(a,m): m\mid 420,\ m\ge2,\ 0\le a<m\}\). Then
\(|V|=\sigma(420)-1=1344-1=1343\) (the script prints
\texttt{vertici=1343}). For each \(k\in\{3,\dots,8\}\) let \(G_k\) be
the graph on \(V\) with an edge between \((a,m)\) and \((a',m')\) iff
\[\gcd(m,m')\le k-1\quad\text{and}\quad \gcd(m,m')\nmid (a-a').\] By
Corollary 2, \textbf{a counterexample of size \(k\) yields a
\(k\)-clique in \(G_k\)}; so if \(G_k\) has no \(k\)-clique, the Theorem
holds for that \(k\). Nothing outside \(V\) needs to be examined: Lemma
1 maps every counterexample into \(V\).

\textbf{Algorithm.} Vertices are indexed \(0,\dots,1342\) in a fixed
order (moduli increasing, residues increasing). \texttt{adj{[}i{]}} is
the bitset of indices \(j>i\) adjacent to \(i\) in \(G_k\), computed
directly from the definition with exact integer arithmetic
(\texttt{math.gcd}, \texttt{\%}). The recursion
\texttt{ricorsione(parziale,\ candidati)} starts with the empty tuple
and the candidate set \(=V\); it iterates over candidates \(j\) in
increasing order and recurses with \texttt{parziale+{[}j{]}} and
\texttt{candidati\ \&\ adj{[}j{]}}.

\emph{Correctness (every \(k\)-clique is found).} By induction on the
length of \texttt{parziale}: the invariant is that \texttt{candidati}
equals the set of indices greater than the last element of
\texttt{parziale} (all of \(V\) if it is empty) that are adjacent to
every element of \texttt{parziale}. Initially true. If it holds for
\texttt{parziale} and \(j\) is a candidate, then
\texttt{candidati\ \&\ adj{[}j{]}} is the set of indices that are \(>j\)
(since \texttt{adj{[}j{]}} only contains indices \(>j\)), adjacent to
every element of \texttt{parziale} (since they are in
\texttt{candidati}) and adjacent to \(j\); this is the invariant for
\texttt{parziale+{[}j{]}}. Therefore every increasing tuple
\((v_1<\dots<v_k)\) of pairwise adjacent vertices is visited (each
\(v_{t+1}\) lies in the candidate set of the prefix
\((v_1,\dots,v_t)\)), and it is visited exactly once (a tuple is reached
only through its unique chain of prefixes). Every \(k\)-clique of
\(G_k\) is exactly one such tuple. Whenever \texttt{len(parziale)==k}
the tuple is recorded as a survivor.

\textbf{Pruning rules (certificate item 2).} None beyond the definition:
a partial tuple is extended only by vertices adjacent (in \(G_k\)) to
all its members. A set containing a non-adjacent pair contains either
two classes that meet or a pair with gcd \(\ge k\), so it is not (a
reduction of) a counterexample; hence nothing that can be completed to a
counterexample is discarded. No symmetry reduction, no ordering
heuristic beyond the fixed vertex order, no extra rules.

\textbf{Node counts and wall clock (certificate item 3).} A node is one
recursive call with nonempty \texttt{parziale}, i.e.~the number of
nonempty pairwise-adjacent increasing tuples of length \(\le k\) (so it
equals \(\sum_{t=1}^{k}\#\{t\text{-cliques of }G_k\}\); e.g.~for
\(k=3\): \(7118 = 1343 + 5775 + 0\), i.e.~\(G_3\) has 5775 edges and no
triangle). Run 1 (\texttt{run1.txt}), Apple M4, Python 3.14.7, exact
integer arithmetic:

{\def\LTcaptype{none} % do not increment counter
\begin{longtable}[]{@{}llllll@{}}
\toprule\noalign{}
\(k\) & gcd threshold \(\le k-1\) & vertices & nodes & survivors &
time \\
\midrule\noalign{}
\endhead
\bottomrule\noalign{}
\endlastfoot
3 & 2 & 1343 & 7118 & 0 & 0.10 s \\
4 & 3 & 1343 & 20241 & 0 & 0.11 s \\
5 & 4 & 1343 & 56208 & 0 & 0.12 s \\
6 & 5 & 1343 & 234719 & 0 & 0.20 s \\
7 & 6 & 1343 & 806717 & 0 & 0.46 s \\
8 & 7 & 1343 & 3523544 & 0 & 1.51 s \\
\end{longtable}
}

Total wall clock of the whole certificate run: 2.5 s (far below the
10-minute limit). Command:
\texttt{.venv/bin/python\ runs/p4\_c3/sandbox/ricerca\_clique\_420.py\ certificato}.

\textbf{Rerun (certificate item 4).} A second run (\texttt{run2.txt})
reproduced every node count and survivor count exactly (verified by
\texttt{diff} after stripping the timing field: \texttt{IDENTICI});
times were 0.10/0.11/0.12/0.20/0.43/1.50 s.

\textbf{Survivors (certificate item 5).} None at any
\(k\in\{3,\dots,8\}\), so there is nothing to decide.

\textbf{Conclusion of Step 2.} For each \(k\in\{3,\dots,8\}\), \(G_k\)
has no \(k\)-clique; by Corollary 2 there is no counterexample of size
\(k\), i.e.~the Theorem holds for \(k=3,4,5,6,7,8\). \(\blacksquare\)

\subparagraph{Step 3. Non-vacuity of the machinery (not required for C3,
included as evidence the code is not silently
broken)}\label{step-3.-non-vacuity-of-the-machinery-not-required-for-c3-included-as-evidence-the-code-is-not-silently-broken}

Running the \emph{same, unmodified} enumeration with the gcd threshold
relaxed from \(k-1\) to \(k\) (mode \texttt{non\_vacuita}), the trivial
partition \(0,1,\dots,k-1 \pmod k\) --- which is an admissible
configuration with all gcds equal to \(k\) --- is found among the
cliques: \(k=3\): 6688 cliques, \(k=4\): 7917, \(k=5\): 17856, \(k=6\):
17329, \(k=7\): 34056 (universe \(N=420\)), each run printing
\texttt{partizione\ banale\ trovata:\ True}. For \(k=8\) the modulus 8
does not divide 420, so the universe was enlarged to the divisors of 840
(2879 vertices): 32 097 759 nodes, 58 425 cliques, 14.7 s, trivial
partition found. (The certificate of Step 2 uses only \(N=420\), which
Lemma 1 justifies for gcds \(\le7\).)

\subparagraph{Step 4. Independent second implementation (cross-check,
not part of the
certificate)}\label{step-4.-independent-second-implementation-cross-check-not-part-of-the-certificate}

\texttt{ricerca\_moduli\_residui.py} uses a different algorithm: it
first enumerates non-decreasing \(k\)-tuples of moduli (divisors of 420,
\(\ge2\)) with all pairwise gcds in \([2,\text{threshold}]\), then
solves the residue CSP \(a_i\not\equiv a_j \pmod{\gcd(m_i,m_j)}\) by
backtracking with the translation normalisation \(a_1=0\). For
thresholds \(k-1\) it examined 51, 77, 221, 293, 1249, 1668 modulus
tuples for \(k=3,\dots,8\) and found 0 families in every case (2.9 s
total). With threshold \(k\) it finds e.g.~\(((0,3),(1,3),(2,3))\),
\(((0,5),\dots,(4,5))\), \(((0,7),\dots,(6,7))\). Its counts are not
directly comparable with those of the clique search (different
normalisation); only the zero/non-zero outcomes are compared, and they
agree.

\subparagraph{\texorpdfstring{Step 5. Hand proofs for \(k=3\) and
\(k=4\) (so the cell does not depend on cells 1--2 being
accepted)}{Step 5. Hand proofs for k=3 and k=4 (so the cell does not depend on cells 1--2 being accepted)}}\label{step-5.-hand-proofs-for-k3-and-k4-so-the-cell-does-not-depend-on-cells-12-being-accepted}

\textbf{\(k=3\).} Suppose three pairwise disjoint classes have all
\(g_{ij}\le 2\). By (P1)-type reasoning (\(g_{ij}=1\) would force the
classes to meet), \(g_{ij}=2\) for all pairs, so \(g_{ij}\nmid a_i-a_j\)
says \(a_i-a_j\) is odd for all \(i<j\). Then \(a_1,a_2,a_3\) have
pairwise different parities, impossible for three integers with only two
parities. Hence some \(g_{ij}\ge3\).

\textbf{\(k=4\).} Suppose four pairwise disjoint classes have all
\(g_{ij}\in\{2,3\}\) (values \(\le3\) and \(\ne1\)). Let
\(E=\{i: 2\mid m_i\}\), \(T=\{i:3\mid m_i\}\). (a) If \(g_{ij}=2\) then
\(i,j\in E\); if \(g_{ij}=3\) then \(i,j\in T\). Hence every pair lies
in \(E\) or in \(T\). (b) If \(i\ne j\) both lie in \(E\cap T\) then
\(6\mid g_{ij}\), contradicting \(g_{ij}\le3\); so \(|E\cap T|\le1\).
(c) Two indices \(i,j\in E\) have \(2\mid g_{ij}\), so \(g_{ij}=2\) and
\(a_i\not\equiv a_j\pmod 2\); thus residues mod 2 of \(E\) are distinct
and \(|E|\le2\). Similarly for \(i,j\in T\), \(3\mid g_{ij}\) gives
\(g_{ij}=3\) and \(a_i\not\equiv a_j \pmod 3\), so \(|T|\le3\). (d) By
(a) no \(i\in E\setminus T\) and \(j\in T\setminus E\) can coexist (the
pair would lie in neither \(E\) nor \(T\)). Hence \(E\subseteq T\) or
\(T\subseteq E\); with (a) (every index lies in \(E\cup T\), as it forms
a pair with some other index) we get \(\{1,2,3,4\}=T\) or
\(\{1,2,3,4\}=E\), contradicting \(|T|\le3\), \(|E|\le2\). Hence some
\(g_{ij}\ge4\).

(These two cases are also covered by the search at \(k=3,4\).)

\subparagraph{Scope and honesty}\label{scope-and-honesty}

Claimed: the Theorem for \(k=3,4,5,6,7,8\) (cell 3 in full, including
the sizes of cells 1--2). Certified by: Lemma 1 (complete proof above) +
exhaustive enumeration with the certificate items 1--5 + rerun. Nothing
is claimed for \(k\ge9\); the reduction to \(N=420\) is valid only for
gcds \(\le7\), and for larger \(k\) one would use
\(N=\mathrm{lcm}(1,\dots,k-1)\) in Lemma 1 (same proof).

\subparagraph{3. Verification: instructions, dependencies,
timings}\label{verification-instructions-dependencies-timings}

Python 3 standard library only. Scripts (also saved in
\texttt{runs/p4\_c3/sandbox/} and re-run in
\texttt{runs/p4\_c3/verifica/}): - code\_1 (python, rigor
\texttt{exact}): Implementation 1 (certificate): exhaustive enumeration
of all k-cliques in G\_k on the 1343 classes (a,m), m\textbar420,
m\textgreater=2; no pruning beyond the edge definition; prints node
counts, survivors, wall clock; mode non\_vacuita relaxes the gcd
threshold to k and checks that the trivial partition 0..k-1 (mod k) is
found. File runs/p4\_c3/sandbox/ricerca\_clique\_420.py; outputs
run1.txt, run2.txt, non\_vacuita.txt, non\_vacuita\_k8\_N840.txt. -
code\_2 (python, rigor \texttt{exact}): Implementation 2 (independent
cross-check, different method): enumerate non-decreasing k-tuples of
moduli dividing 420 with pairwise gcd in {[}2, threshold{]}, then solve
the residue CSP a\_i != a\_j (mod gcd) by backtracking with a\_1 = 0.
Found 0 families at thresholds k-1 for k=3..8 (51, 77, 221, 293, 1249,
1668 modulus tuples examined; 2.9 s total) and finds the trivial
partitions at threshold k. File
runs/p4\_c3/sandbox/ricerca\_moduli\_residui.py; outputs
impl2\_run1.txt, impl2\_non\_vacuita.txt.

Trusted re-runs by the orchestrator (exit code, wall clock, output): -
orchestrator re-ran code\_1.py (python3, clean copy of the researcher
sandbox): exit 0 in 2.5s; stdout: 'k=3 soglia\_gcd\textless=2
vertici=1343 nodi=7118 sopravvissuti=0 tempo=0.10s\nk=4
soglia\_gcd\textless=3 vertici=1343 nodi=20241 sopravvissuti=0
tempo=0.10s\nk=5 soglia\_gcd\textless=4 vertici=1343 nodi=56208
sopravvissuti - orchestrator re-ran code\_2.py (python3, clean copy of
the researcher sandbox): exit 0 in 3.0s; stdout: 'k=3
soglia\_gcd\textless=2 multinsiemi\_moduli=51 famiglie(a1=0)=0
tempo=0.00s\nk=4 soglia\_gcd\textless=3 multinsiemi\_moduli=77
famiglie(a1=0)=0 tempo=0.00s\nk=5 soglia\_gcd\textless=4
multinsiemi\_moduli=221 famiglie(a1=0)=

\subparagraph{4. Sources and
contribution}\label{sources-and-contribution}

\begin{itemize}
\tightlist
\item
  Problem statement of Problema 4 (criterion (*) and definitions), file
  problema-4/enunciato.md
\item
  arXiv:2607.24655 (Fornal--Sun) --- abstract only, used for
  literature\_position, no result relied upon
\item
  arXiv:2608.15873 (Menon) --- abstract only, context for the group
  form, no result relied upon
\item
  Position with respect to the literature: The listed abstracts do not
  settle this cell in a usable way: {[}2607.24655{]} (Fornal--Sun)
  proves the asymptotic bound max gcd ≫ k·exp(−(2+o(1))√(log k/log log
  k)), which is relevant to C6(b), not to small k; {[}2608.15873{]}
  (Menon) concerns the group form (C6(c)) and the tuple (6,6,6,10,15);
  {[}2603.26043{]} and {[}1511.04293{]} concern disjoint covering
  systems with one repeated modulus, unrelated to the gcd question. I
  did not read the full papers and cite nothing from them; my approach
  (finite reduction + exhaustive certified search) follows the column's
  own hand-in rules rather than any of these papers. I recall that the k
  ≤ 8 case is probably classical (Huhn--Megyesi / Z.-W. Sun on disjoint
  residue classes) but I have not verified any such source, so I do not
  cite it as read and I rely on it for nothing.
\item
  Contribution: the proof above is written out in full by the team's
  Researcher and checked by two independent judges and by a human.
\end{itemize}

\subparagraph{5. Limits and unresolved
parts}\label{limits-and-unresolved-parts}

Gaps declared by the author (all accepted by the judges): - The search
is only as exhaustive as the code is correct; I proved the enumeration
invariant in prose (Step 2) but a hostile reader must still read the
\textasciitilde80 lines of Python. The independent second implementation
and the non-vacuity runs mitigate but do not replace this. - The k=8
non-vacuity check needed a larger universe (divisors of 840) because 8
does not divide 420; this is expected from Lemma 1 (valid only for gcds
≤ 7) and does not affect the certificate, but it is a deviation from
`run the machinery unmodified' in the strict C5 sense (not required for
C3). - The survivor counts of the two implementations in non-vacuity
mode are not directly comparable (unordered cliques
vs.~a\_1=0-normalised ordered families with repeated moduli), so the
cross-check compares only zero/non-zero outcomes at thresholds k−1, plus
presence of the trivial partition. - Literature: I did not fetch or read
the full papers listed; I recall that the small-k cases are probably
classical (Huhn--Megyesi / Z.-W. Sun) but have not verified this and
rely on it for nothing.

\subsubsection{Decisioni degli agenti}
\paragraph{Run \texttt{p4\_c3}}
\begin{description}
\item[Iterazione 1 — attempt\_001]
\textbf{Researcher}: famiglia \texttt{exact\_computation}, sotto-obiettivo: Prove the statement for every k with 3 ≤ k ≤ 8 (self-contained, not relying on cells 1–2 being verified): reduction lemma to moduli dividing 420 with pairwise gcds preserved, then a fully certified exhaustive clique search (finite set, proof of exhaustiveness, no pruning beyond the definition, node counts, wall clock, rerun, zero survivors), plus hand proofs for k=3,4 and a non-vacuity test.. Stato dichiarato: \texttt{CELL\_SOLVED\_CANDIDATE}.
\emph{Perché questa scelta}: First natural subgoal (no blocker recorded, no verified claims). The hand-in rules of this column explicitly admit a computation with an exhaustiveness certificate; the only genuinely mathematical step is the finite reduction, which I prove in full. I deliberately used NO pruning beyond the definition of a counterexample (so rule 2 of the certificate is trivially satisfied) and no symmetry reduction, because the search is small (3.5 million nodes at k=8). Since highest\_verified\_cell = 0, I do not assume cells 1–2 and give short hand proofs for k=3,4 as well; the computation covers them independently.
\emph{Rispetto allo stato dell'arte}: The listed abstracts do not settle this cell in a usable way: [2607.24655] (Fornal–Sun) proves the asymptotic bound max gcd ≫ k·exp(−(2+o(1))√(log k/log log k)), which is relevant to C6(b), not to small k; [2608.15873] (Menon) concerns the group form (C6(c)) and the tuple (6,6,6,10,15); [2603.26043] and [1511.04293] concern disjoint covering systems with one repeated modulus, unrelated to the gcd question. I did not read the full papers and cite nothing from them; my approach (finite reduction + exhaustive certified search) follows the column's own hand-in rules rather than any of these papers. I recall that the k ≤ 8 case is probably classical (Huhn–Megyesi / Z.-W. Sun on disjoint residue classes) but I have not verified any such source, so I do not cite it as read and I rely on it for nothing.
\textbf{Referee}: verdetto \texttt{UNKNOWN\_STATUS} (review\_status \texttt{READY\_FOR\_HUMAN}).
\emph{Prossimo passo richiesto}: Human reviews the exact target, proof and evidence, then approves explicit claims
\end{description}

\subsubsection{Codice su cui poggia il tentativo}
\begin{lstlisting}[caption={attempt\_001 code\_1 [python, rigore exact] — Implementation 1 (certificate): exhaustive enumeration of all k-cliques in G\_k on the 1343 — approvato dal Referee}]
"""
Ricerca esaustiva per il problema 4, parte 3 (k <= 8).

COSA: enumera tutte le famiglie di k classi a_i (mod m_i), a due a due disgiunte,
con m_i | 420, 2 <= m_i, 0 <= a_i < m_i, e con gcd(m_i, m_j) <= soglia per ogni coppia.
PERCHE': il Lemma di riduzione (nella prova) mostra che ogni controesempio di taglia k
(famiglia disgiunta con tutti i gcd <= k-1 <= 7) si riduce a una tale famiglia con
m_i | 420 = 2^2 * 3 * 5 * 7. Se non ne esiste nessuna, la tesi vale per quel k.

Metodo: grafo G sui vertici (a, m); arco fra due vertici se e solo se le due classi
sono disgiunte E il gcd dei moduli e' <= soglia. Enumerazione di TUTTE le clique di
taglia k come tuple crescenti di indici (nessuna potatura oltre la definizione di arco).
Aritmetica: esatta (interi Python).
"""
import sys
import time
from math import gcd

N_RIDUZIONE = 420


def divisori(n):
    """Tutti i divisori di n, in ordine crescente."""
    return [d for d in range(1, n + 1) if n % d == 0]


def costruisci_vertici(n):
    """Lista delle classi (a, m) con m | n, m >= 2, 0 <= a < m."""
    return [(a, m) for m in divisori(n) if m >= 2 for a in range(m)]


def compatibili(v, w, soglia):
    """Vero se le classi v, w sono disgiunte e gcd dei moduli <= soglia (criterio (*))."""
    (a, m), (b, n) = v, w
    g = gcd(m, n)
    return g <= soglia and (a - b) % g != 0


def costruisci_adiacenza(vertici, soglia):
    """Bitset di adiacenza: bit j di adj[i] acceso se i e j compatibili (solo j > i)."""
    adj = []
    for i, v in enumerate(vertici):
        bits = 0
        for j in range(i + 1, len(vertici)):
            if compatibili(v, vertici[j], soglia):
                bits |= 1 << j
        adj.append(bits)
    return adj


def enumera_clique(adj, k, raccogli):
    """
    Enumera tutte le k-clique come tuple crescenti di indici.
    Ritorna (numero_nodi_visitati, lista_clique). Un nodo = una chiamata ricorsiva
    con clique parziale non vuota (ogni prefisso visitato una sola volta).
    """
    nodi = 0
    trovate = []

    def ricorsione(parziale, candidati):
        nonlocal nodi
        nodi += 1
        if len(parziale) == k:
            trovate.append(tuple(parziale))
            return
        c = candidati
        while c:
            j = (c & -c).bit_length() - 1
            c &= c - 1
            ricorsione(parziale + [j], candidati & adj[j])

    tutti = (1 << len(adj)) - 1
    ricorsione([], tutti)
    return nodi - 1, trovate  # il nodo radice (clique vuota) non si conta


def esegui(k, soglia, vertici, stampa_sopravvissuti):
    """Una taglia k con una soglia gcd; stampa conteggi, tempo e sopravvissuti."""
    t0 = time.time()
    adj = costruisci_adiacenza(vertici, soglia)
    nodi, clique = enumera_clique(adj, k, True)
    dt = time.time() - t0
    print(f"k={k} soglia_gcd<={soglia} vertici={len(vertici)} nodi={nodi} "
          f"sopravvissuti={len(clique)} tempo={dt:.2f}s")
    if stampa_sopravvissuti:
        for cl in clique:
            print("   ", [vertici[i] for i in cl])
    return nodi, clique


def main():
    modo = sys.argv[1] if len(sys.argv) > 1 else "certificato"
    # Universo: divisori di N (default 420). Per il test di non vacuita' a k=8
    # serve N=840, perche' 8 non divide 420 (il certificato usa sempre 420).
    universo = int(sys.argv[3]) if len(sys.argv) > 3 else N_RIDUZIONE
    vertici = costruisci_vertici(universo)
    if modo == "certificato":
        # Controesempio di taglia k: tutti i gcd <= k-1.
        for k in range(3, 9):
            esegui(k, k - 1, vertici, stampa_sopravvissuti=True)
    elif modo == "non_vacuita":
        # Soglia gcd <= k: la partizione 0..k-1 (mod k) DEVE comparire.
        k = int(sys.argv[2])
        nodi, clique = esegui(k, k, vertici, stampa_sopravvissuti=False)
        atteso = tuple(vertici.index((a, k)) for a in range(k))
        print("   partizione banale trovata:", atteso in set(clique))
    else:
        raise SystemExit("modo: certificato | non_vacuita K [N]")


if __name__ == "__main__":
    main()

\end{lstlisting}
\noindent\emph{Riesecuzione fidata dell'orchestratore}: \texttt{orchestrator re-ran code\_1.py (python3, clean copy of the researcher sandbox): exit 0 in 2.5s; stdout: 'k=3 soglia\_gcd\textless{}=2 vertici=1343 nodi=7118 sopravvissuti=0 tempo=0.10s\textbackslash{}nk=4 soglia\_gcd\textless{}=3 vertici=1343 nodi=20241 sopravvissuti=0 tempo=0.10s\textbackslash{}nk=5 soglia\_gcd\textless{}=4 vertici=1343 nodi=56208 sopravvissuti=0 tempo=0.12s\textbackslash{}nk=6 soglia\_gcd\textless{}=5 vertici=1343 nodi=234719 sopravvissuti=0 tempo=0.20s\textbackslash{}nk=7 soglia\_g}
\begin{lstlisting}[caption={attempt\_001 code\_2 [python, rigore exact] — Implementation 2 (independent cross-check, different method): enumerate non-decreasing k-t — approvato dal Referee}]
"""
Seconda implementazione (metodo diverso) per il problema 4, parte 3.

COSA: cerca famiglie disgiunte di k classi con m_i | 420, tutti i gcd(m_i,m_j) <= soglia.
COME: (1) enumera i multinsiemi non decrescenti di moduli (divisori di 420, >= 2) con
gcd a coppie in [2, soglia]; (2) per ciascuno risolve il CSP sui residui
a_i (mod m_i) con vincolo a_i != a_j (mod gcd), normalizzando a_1 = 0 (traslazione).
PERCHE': controllo incrociato del risultato dell'enumerazione di clique, con un
algoritmo strutturalmente diverso. Aritmetica esatta.
"""
import sys
import time
from math import gcd

MODULI = [d for d in range(2, 421) if 420 % d == 0]


def multinsiemi_moduli(k, soglia):
    """Tuple non decrescenti di k moduli con gcd a coppie in [2, soglia]."""
    def estendi(parziale, da):
        if len(parziale) == k:
            yield tuple(parziale)
            return
        for m in MODULI:
            if m < da:
                continue
            if all(2 <= gcd(m, q) <= soglia for q in parziale):
                yield from estendi(parziale + [m], m)
    yield from estendi([], 2)


def residui_compatibili(moduli):
    """Tutte le assegnazioni di residui (a_1 = 0) che rendono le classi disgiunte."""
    soluzioni = []

    def estendi(res):
        i = len(res)
        if i == len(moduli):
            soluzioni.append(tuple(res))
            return
        for a in range(moduli[i]) if i > 0 else [0]:
            if all((a - res[j]) % gcd(moduli[i], moduli[j]) != 0 for j in range(i)):
                estendi(res + [a])
    estendi([])
    return soluzioni


def esegui(k, soglia):
    """Conta multinsiemi di moduli esaminati e famiglie trovate."""
    t0 = time.time()
    n_moduli, famiglie = 0, []
    for moduli in multinsiemi_moduli(k, soglia):
        n_moduli += 1
        for res in residui_compatibili(moduli):
            famiglie.append(tuple(zip(res, moduli)))
    dt = time.time() - t0
    print(f"k={k} soglia_gcd<={soglia} multinsiemi_moduli={n_moduli} "
          f"famiglie(a1=0)={len(famiglie)} tempo={dt:.2f}s")
    return famiglie


if __name__ == "__main__":
    if len(sys.argv) > 1:
        k, soglia = int(sys.argv[1]), int(sys.argv[2])
        fam = esegui(k, soglia)
        print("   esempio:", fam[0] if fam else None)
    else:
        for k in range(3, 9):
            esegui(k, k - 1)

\end{lstlisting}
\noindent\emph{Riesecuzione fidata dell'orchestratore}: \texttt{orchestrator re-ran code\_2.py (python3, clean copy of the researcher sandbox): exit 0 in 3.0s; stdout: 'k=3 soglia\_gcd\textless{}=2 multinsiemi\_moduli=51 famiglie(a1=0)=0 tempo=0.00s\textbackslash{}nk=4 soglia\_gcd\textless{}=3 multinsiemi\_moduli=77 famiglie(a1=0)=0 tempo=0.00s\textbackslash{}nk=5 soglia\_gcd\textless{}=4 multinsiemi\_moduli=221 famiglie(a1=0)=0 tempo=0.02s\textbackslash{}nk=6 soglia\_gcd\textless{}=5 multinsiemi\_moduli=293 famiglie(a1=0)=0 tempo=0.06s\textbackslash{}nk=7 soglia\_gcd}

\section{arXiv literature consulted}
\begin{thebibliography}{99}
\bibitem{2603.26043v1} Yu Hashimoto. \emph{Finiteness of Disjoint Covering Systems with Precisely One Repeated Modulus}. arXiv:2603.26043v1 (2026). \url{http://arxiv.org/abs/2603.26043v1}
\bibitem{1511.04293v1} Shalosh B. Ekhad, Aviezri S. Fraenkel, Doron Zeilberger. \emph{Searching for Disjoint Covering Systems with Precisely One Repeated Modulus}. arXiv:1511.04293v1 (2015). \url{http://arxiv.org/abs/1511.04293v1}
\bibitem{2607.24655v1} Jan Fornal, Yu-Chen Sun. \emph{On the problem of large gcd for disjoint residue classes}. arXiv:2607.24655v1 (2026). \url{http://arxiv.org/abs/2607.24655v1}
\bibitem{2608.15873v1} Murali Menon. \emph{Two Questions on \$G\$-harmonic Tuples}. arXiv:2608.15873v1 (2026). \url{http://arxiv.org/abs/2608.15873v1}
\end{thebibliography}
\end{document}
